\chapter{The two-purity criterion}
\label{chap:criterion}

The doubling step from seed purities to a uniform spectral gap depends only on algebraic relations
among partition functions, not on the microscopic Hilbert space of a spin chain. Formalized in
\texttt{CertifiedGaps/PurityCriterion.lean}, the criterion applies to any family of finite spectra
\(E_L \colon \iota_L \to \R\), \(L \in \N\), over finite nonempty index types whose partition
functions satisfy the three spatial moment inequalities of Definition~\ref{def:moment-bounds}.

\section{Partition functions, purities and gaps}

For each system size, the partition function sums the Boltzmann weights of the energy levels, and
the spatial and thermal purities compare partition functions at doubled length or doubled inverse
temperature. To state spectral gap bounds without sorting eigenvalues, the predicate
\(\GapGT(e, c)\) asks for a single ground state separated from all other levels by more than
\(c\).

\begin{definition}[Partition function]
  \label{def:partition-function}
  \lean{CertifiedGaps.partitionFunction}
  \leanok
  The partition function of the spectrum \(E_L\) is
  \(Z_L(\beta) = \sum_{m \in \iota_L} e^{-\beta E_L(m)}\).
\end{definition}

\begin{definition}[Purities]
  \label{def:purities}
  \lean{CertifiedGaps.thermalPurity, CertifiedGaps.spatialPurity}
  \leanok
  \uses{def:partition-function}
  The thermal purity is \(\Pt(L, \beta) = Z_L(2\beta)/Z_L(\beta)^2\) and the spatial purity is
  \(\Ps(n, \beta) = Z_{2n}(\beta)/Z_n(\beta)^2\). Their complements \(1 - \Pt\) and \(1 - \Ps\) are
  the \emph{thermal defect} and \emph{spatial defect}.
\end{definition}

\begin{definition}[Gap predicate]
  \label{def:gapgt}
  \lean{CertifiedGaps.GapGT}
  \leanok
  For \(e \colon \kappa \to \R\) and \(c \in \R\), the predicate \(\GapGT(e, c)\) holds if there is
  an index \(m_0\) with \(e(m_0) + c < e(m)\) for every \(m \ne m_0\).
\end{definition}

\begin{lemma}
  \label{lem:gapgt-monotone}
  \lean{CertifiedGaps.GapGT.lt_sub_of_monotone}
  \leanok
  \uses{def:gapgt}
  If \(e_0 \le e_1 \le \cdots \le e_{N+1}\) is ordered and \(\GapGT(e, c)\) holds, then
  \(e_1 - e_0 > c\).
\end{lemma}
\begin{proof}
  \leanok
  If the minimizing index is \(m_0 = 0\), evaluating the defining inequality at \(m = 1\) gives the
  bound. Otherwise, evaluating it at \(m = 0\) and combining it with \(e_0 \le e_{m_0}\) and
  \(e_0 \le e_1\) again gives the bound.
\end{proof}

\begin{lemma}[Elementary bounds]
  \label{lem:purity-basic}
  \lean{CertifiedGaps.partitionFunction_pos, CertifiedGaps.thermalPurity_pos,
    CertifiedGaps.spatialPurity_pos, CertifiedGaps.thermalPurity_le_one}
  \leanok
  \uses{def:purities, lem:power-sums}
  For every \(L\) and \(\beta\), \(Z_L(\beta) > 0\), both purities are positive, and
  \(\Pt(L, \beta) \le 1\).
\end{lemma}
\begin{proof}
  \leanok
  Positivity holds because each Boltzmann weight \(e^{-\beta E_L(m)}\) is positive and the index
  type is nonempty. Since \(Z_L(2\beta) = \sum_m (e^{-\beta E_L(m)})^2\), the bound \(\Pt \le 1\)
  reduces to the inequality \(\sum p_m^2 \le (\sum p_m)^2\) for \(p_m \ge 0\).
\end{proof}

\section{The squaring lemma}

When a positive operator with trace one and eigenvalues \(p_i\) is squared and renormalized, its
eigenvalues become \(p_i^2/P\) and its purity defect \(1 - P\) is bounded after one step by the
rational map \(f(z) = z^2/[2(1-z)^2]\). Stated in homogeneous form for the power sums
\(S_1, S_2, S_4\), the same inequality applies to both finite Boltzmann sums and unnormalized
spatial partition functions.

\begin{lemma}[Power sums]
  \label{lem:power-sums}
  \lean{CertifiedGaps.sq_sub_sum_sq_finset, CertifiedGaps.sq_sub_sum_sq_of_hasSum}
  \leanok
  For a nonnegative summable family \(p_i\) with power sums \(S_k = \sum_i p_i^k\),
  \[
    S_2 \le S_1^2 \qquad\text{and}\qquad 2\,(S_2^2 - S_4) \le (S_1^2 - S_2)^2.
  \]
\end{lemma}
\begin{proof}
  \leanok
  For finite families, induct on the support. Adding a term \(a \ge 0\) to a family with sum \(s\)
  and difference \(T = s^2 - S_2 \ge 0\) increases the right-hand side by at least as much as the
  left-hand side, since their difference changes by nonnegative multiples of \(a s T\) and
  \(a^2 T\). The summable case follows by taking the limit over finite subsets, because both
  inequalities are closed.
\end{proof}

\begin{definition}[Squaring bound]
  \label{def:squaring-defect}
  \lean{CertifiedGaps.squaringDefect}
  \leanok
  The squaring-defect function is \(f(z) = z^2 / \bigl(2 (1 - z)^2\bigr)\).
\end{definition}

\begin{lemma}
  \label{lem:f-mono}
  \lean{CertifiedGaps.squaringDefect_mono}
  \leanok
  \uses{def:squaring-defect}
  The function \(f\) is nondecreasing on \([0, 1)\).
\end{lemma}
\begin{proof}
  \leanok
  For \(0 \le z \le w < 1\), expanding the products shows \(z (1 - w) \le w (1 - z)\); squaring
  both sides and dividing by the positive denominators gives the bound.
\end{proof}

\begin{lemma}[Squaring, homogeneous form]
  \label{lem:squaring-homogeneous}
  \lean{CertifiedGaps.one_sub_div_le_squaringDefect}
  \leanok
  \uses{def:squaring-defect}
  If \(s_1, s_2 > 0\) and \(2 (s_2^2 - s_4) \le (s_1^2 - s_2)^2\), then
  \(1 - s_4/s_2^2 \le f(1 - s_2/s_1^2)\).
\end{lemma}
\begin{proof}
  \leanok
  Rewriting the left- and right-hand sides as \((s_2^2 - s_4)/s_2^2\) and
  \((s_1^2 - s_2)^2 / (2 s_2^2)\) reduces the claim to the assumed inequality divided by the
  positive denominator.
\end{proof}

\begin{lemma}[Squaring]
  \label{lem:squaring}
  \lean{CertifiedGaps.one_sub_squaredPurity_le_squaringDefect}
  \leanok
  \uses{lem:power-sums, lem:squaring-homogeneous}
  Let \(p_i \ge 0\) with \(\sum_i p_i = 1\) and purity \(P = \sum_i p_i^2\). The squared
  distribution \(p_i^2 / P\) has purity \(P'\) satisfying \(1 - P' \le f(1 - P)\).
\end{lemma}
\begin{proof}
  \leanok
  The families \(p_i^2\) and \(p_i^4\) are summable because \(0 \le p_i \le 1\), and \(P > 0\)
  follows from \(\sum p_i = 1\). Writing the new purity as \(P' = S_4 / S_2^2\) and applying
  Lemma~\ref{lem:power-sums} together with Lemma~\ref{lem:squaring-homogeneous} gives the bound.
\end{proof}

\begin{lemma}
  \label{lem:f-numerical}
  \lean{CertifiedGaps.squaringDefect_le}
  \leanok
  \uses{def:squaring-defect}
  For \(0 \le \varepsilon \le 1/100\), \(f\bigl(1 - (1 - \varepsilon)^3\bigr) \le 5 \varepsilon^2\).
\end{lemma}
\begin{proof}
  \leanok
  Setting \(u = (1 - \varepsilon)^3\) gives \(1 - 3\varepsilon \le u \le 1\), which implies
  \((1 - u)^2 \le 9 \varepsilon^2\) and \(9 \le 10 u^2\); combining the two bounds yields the
  inequality.
\end{proof}

\begin{lemma}[Squaring step]
  \label{lem:squaring-step}
  \lean{CertifiedGaps.squaring_step_of_le, CertifiedGaps.squaring_step}
  \leanok
  \uses{lem:power-sums, lem:squaring-homogeneous, lem:f-mono, lem:f-numerical}
  Let \(s_1, s_2 > 0\) with \(s_2 \le s_1^2\) and \(2 (s_2^2 - s_4) \le (s_1^2 - s_2)^2\) (for
  instance the power sums of a nonnegative summable family), and let
  \(0 \le \varepsilon \le 1/100\). If \(s_2 / s_1^2 \ge (1 - \varepsilon)^3\), then
  \(s_4 / s_2^2 \ge 1 - 5 \varepsilon^2\).
\end{lemma}
\begin{proof}
  \leanok
  By Lemma~\ref{lem:squaring-homogeneous}, the defect after squaring is at most
  \(f(1 - s_2/s_1^2)\). Monotonicity makes this at most \(f(1 - (1-\varepsilon)^3)\) by
  Lemma~\ref{lem:f-mono}, which is bounded by \(5\varepsilon^2\) by Lemma~\ref{lem:f-numerical}.
\end{proof}

\section{The spatial moment interface}

The doubling recurrence uses the spatial transfer operator only through three inequalities among
the partition functions \(Z_L(\beta)\) at a fixed inverse temperature, namely two power-sum
inequalities for spatial squaring and one power-sum monotonicity inequality to interpolate across
intermediate even lengths. Because each inequality involves only finitely many values
\(Z_L(\beta)\), the condition is preserved under pointwise limits
(Lemma~\ref{lem:moment-bounds-limit}), which Chapter~\ref{chap:rings} uses to pass from finite
Trotter approximations to the continuous-time partition function.

\begin{definition}[Moment form]
  \label{def:moment-form}
  \lean{CertifiedGaps.HasMomentForm, CertifiedGaps.HasSpatialMomentForm}
  \leanok
  \uses{def:partition-function}
  A sequence \(Z \colon \N \to \R\) has the \emph{moment form} above \(n_0\) if there exist
  \(\mu_i \ge 0\), \(i \in \N\), with convergent sum \(Z(L) = \sum_i \mu_i^L\) for every even
  \(L \ge n_0\). A spectral family has the spatial moment form if \(L \mapsto Z_L(\beta)\) has the
  moment form for every \(\beta > 0\).
\end{definition}

\begin{definition}[Spatial moment inequalities]
  \label{def:moment-bounds}
  \lean{CertifiedGaps.SpatialMomentBounds, CertifiedGaps.HasSpatialMomentBounds}
  \leanok
  \uses{def:partition-function}
  A sequence \(Z\) satisfies the spatial moment inequalities above \(n_0\) if, for all even
  \(n \ge n_0\) and even \(q \ge p \ge n_0\),
  \begin{enumerate}
    \item \(Z(2n) \le Z(n)^2\);
    \item \(2\bigl(Z(2n)^2 - Z(4n)\bigr) \le \bigl(Z(n)^2 - Z(2n)\bigr)^2\);
    \item \(Z(q) \le Z(p)^{q/p}\).
  \end{enumerate}
  A spectral family satisfies them if \(L \mapsto Z_L(\beta)\) does for every \(\beta > 0\).
\end{definition}

\begin{lemma}[Power-sum monotonicity]
  \label{lem:lp-mono}
  \lean{CertifiedGaps.hasSum_pow_le_rpow}
  \leanok
  For \(\mu_i \ge 0\) and \(0 < p \le q\), \(\sum_i \mu_i^q \le \bigl(\sum_i \mu_i^p\bigr)^{q/p}\).
\end{lemma}
\begin{proof}
  \leanok
  Every term \(\mu_i^p\) is bounded by the total sum \(Z_p = \sum_i \mu_i^p\), so
  \(\mu_i^q = (\mu_i^p)^{q/p} \le Z_p^{q/p - 1} \mu_i^p\). Summing over \(i\) gives the bound.
\end{proof}

\begin{proposition}
  \label{prop:moment-form-bounds}
  \lean{CertifiedGaps.HasMomentForm.spatialMomentBounds,
    CertifiedGaps.HasSpatialMomentForm.hasSpatialMomentBounds}
  \leanok
  \uses{def:moment-form, def:moment-bounds, lem:power-sums, lem:lp-mono}
  For \(n_0 > 0\), the moment form above \(n_0\) implies the spatial moment inequalities above
  \(n_0\).
\end{proposition}
\begin{proof}
  \leanok
  Inequalities (1) and (2) follow from Lemma~\ref{lem:power-sums} applied to the family
  \(\mu_i^n\), whose power sums are \(Z(n), Z(2n), Z(4n)\). Inequality (3) follows from
  Lemma~\ref{lem:lp-mono}.
\end{proof}

\begin{lemma}[Closure under pointwise limits]
  \label{lem:moment-bounds-limit}
  \lean{CertifiedGaps.SpatialMomentBounds.of_tendsto, CertifiedGaps.SpatialMomentBounds.mono,
    CertifiedGaps.SpatialMomentBounds.of_tendsto_hasMomentForm}
  \leanok
  \uses{def:moment-bounds, prop:moment-form-bounds}
  If the sequences \(Z_a\) eventually satisfy the spatial moment inequalities above \(n_0 > 0\)
  along a filter (for instance because they have the moment form), and \(Z_a(L) \to Z(L)\) for
  every even \(L \ge n_0\), then \(Z\) satisfies them as well. Moreover, the inequalities above
  \(n_0\) imply those above any \(n_1 \ge n_0\).
\end{lemma}
\begin{proof}
  \leanok
  Each of the three bounds is a non-strict inequality between continuous functions of finitely many
  values \(Z(L)\), and is therefore preserved by pointwise limits. Continuity of the \(q/p\)-th
  power holds because \(q/p \ge 0\).
\end{proof}

\begin{remark}
  Closure under pointwise limits in Lemma~\ref{lem:moment-bounds-limit} avoids formalizing an
  infinite-dimensional path-space transfer operator, since the limit of the Trotter approximations
  of
  \(Z_L(\beta)\) in Chapter~\ref{chap:rings} is taken directly inside the three real-valued
  inequalities.
\end{remark}

\section{Coupling, squaring and the purity-to-gap bound}

Applying the homogeneous squaring lemma to the Boltzmann weights and to the spatial moment
inequalities bounds the thermal purity at doubled inverse temperature and the spatial purity at
doubled length. To advance length and inverse temperature together, two algebraic coupling
identities relate the purities across doubled scales, and the ratio bound
\(e^{-\beta(E_L(m) - E_L(m_0))} \le (1 - \Pt(L, \beta))/\Pt(L, \beta)\) converts a thermal purity
near one into a lower bound on the spectral gap.

\begin{lemma}[Squaring the Gibbs and spatial spectra]
  \label{lem:purity-squaring}
  \lean{CertifiedGaps.thermalPurity_two_mul_ge, CertifiedGaps.spatialPurity_le_one,
    CertifiedGaps.spatialPurity_two_mul_ge}
  \leanok
  \uses{lem:squaring-step, lem:purity-basic, def:moment-bounds}
  Let \(0 \le \varepsilon \le 1/100\). If \(\Pt(L, \beta) \ge (1-\varepsilon)^3\), then
  \(\Pt(L, 2\beta) \ge 1 - 5 \varepsilon^2\). Under the spatial moment inequalities above \(n_0\),
  for even \(n \ge n_0\) and \(\beta > 0\), \(\Ps(n, \beta) \le 1\), and
  \(\Ps(n, \beta) \ge (1-\varepsilon)^3\) implies \(\Ps(2n, \beta) \ge 1 - 5\varepsilon^2\).
\end{lemma}
\begin{proof}
  \leanok
  For the thermal purity, apply Lemma~\ref{lem:squaring-step} to the Boltzmann weights
  \(e^{-\beta E_L(m)}\), whose power sums are \(Z_L(\beta), Z_L(2\beta), Z_L(4\beta)\). For the
  spatial purity, apply the same lemma using inequalities (1) and (2) of
  Definition~\ref{def:moment-bounds}.
\end{proof}

\begin{lemma}[Coupling]
  \label{lem:coupling}
  \lean{CertifiedGaps.spatialPurity_two_mul, CertifiedGaps.thermalPurity_four_mul}
  \leanok
  \uses{def:purities, lem:purity-basic}
  For all \(n\) and \(\beta\),
  \[
    \Ps(n, 2\beta) = \frac{\Ps(n, \beta)^2\, \Pt(2n, \beta)}{\Pt(n, \beta)^2},\qquad
    \Pt(4n, \beta) = \frac{\Pt(2n, \beta)^2\, \Ps(2n, 2\beta)}{\Ps(2n, \beta)^2}.
  \]
\end{lemma}
\begin{proof}
  \leanok
  Expanding each purity by its definition expresses both sides of the first equation in terms of
  the four positive numbers \(Z_n(\beta),
  Z_n(2\beta), Z_{2n}(\beta), Z_{2n}(2\beta)\), and both sides of the second in terms of
  \(Z_{2n}(\beta), Z_{2n}(2\beta),
  Z_{4n}(\beta), Z_{4n}(2\beta)\).
\end{proof}

\begin{lemma}[From thermal purity to a gap]
  \label{lem:purity-to-gap}
  \lean{CertifiedGaps.exp_neg_mul_sub_le_of_isMin}
  \leanok
  \uses{def:purities, lem:purity-basic}
  Let \(\beta \ge 0\), let \(m_0\) minimize \(E_L\), and let \(m \ne m_0\). Then
  \[
    e^{-\beta (E_L(m) - E_L(m_0))} \le \frac{1 - \Pt(L, \beta)}{\Pt(L, \beta)}.
  \]
\end{lemma}
\begin{proof}
  \leanok
  Write \(a_k = e^{-\beta E_L(k)}\), \(Z = \sum a_k\), and \(S = \sum a_k^2\). Because \(a_{m_0}\)
  is the largest Boltzmann weight, \(S \le a_{m_0} Z\), and because \(m \ne m_0\), the two terms
  satisfy \(a_m + a_{m_0} \le Z\). Multiplying gives
  \(a_m S \le (Z - a_{m_0}) S \le Z^2 a_{m_0} - S a_{m_0}\), and dividing by \(a_{m_0} S\) yields
  the inequality.
\end{proof}

\section{Doubling and the uniform gap}

Alternating the two coupling identities with spatial and thermal squaring advances \((n, \beta)\)
to \((2n, 2\beta)\) while mapping a defect bound \(\varepsilon \le 1/100\) to \(5\varepsilon^2\).
Iterating this step produces doubly exponential defect decay along dyadic scales
\((n_0 2^j, \beta_0 2^j)\), and interpolating across each dyadic window \([n_j, 2n_j]\) yields a
uniform gap for all even \(L \ge n_0\).

\begin{definition}[Doubling coefficient]
  \label{def:doubling-coefficient}
  \lean{CertifiedGaps.doublingCoefficient, CertifiedGaps.squaringDefect_couplingDefect_self,
    CertifiedGaps.doublingCoefficient_mono}
  \leanok
  \uses{def:squaring-defect}
  For \(u < 1\) let
  \(G(u) = \tfrac12\bigl((1-u)^{-1} + (1-u)^{-2} + (1-u)^{-3}\bigr)^2\). Then
  \(f\bigl(1 - (1-u)^3\bigr) = G(u)\,u^2\), and \(G\) is increasing.
\end{definition}

\begin{lemma}[Doubling with a coefficient]
  \label{lem:doubling-coefficient}
  \lean{CertifiedGaps.doubling_of_coefficient}
  \leanok
  \uses{def:doubling-coefficient, lem:doubling-step}
  Assume the spatial moment inequalities above \(n_0\). Let \(n \ge n_0\) be even, \(\beta > 0\),
  \(0 \le u < 1\), \(G(u) \le C\) and \(Cu \le 1\). If \(\Ps(n, \beta) \ge 1 - u\) and
  \(\Pt(2n, \beta) \ge 1 - u\), then \(\Ps(2n, 2\beta) \ge 1 - Cu^2\) and
  \(\Pt(4n, 2\beta) \ge 1 - Cu^2\).
\end{lemma}
\begin{proof}
  \leanok
  Apply the exact doubling step (Lemma~\ref{lem:doubling-step}) with \(s = t = u\) and
  \(s' = t' = Cu^2\): the first output is \(f(1 - (1-u)^3) = G(u) u^2 \le Cu^2\), and since
  \(Cu^2 \le u\) the second input is at most the first, so \(f\) is at most \(Cu^2\) there too.
\end{proof}

\begin{lemma}[Doubling]
  \label{lem:doubling}
  \lean{CertifiedGaps.doubling}
  \leanok
  \uses{lem:doubling-coefficient}
  Assume the spatial moment inequalities above \(n_0\). Let \(n \ge n_0\) be even, \(\beta > 0\),
  and \(0 \le \varepsilon \le 1/100\). If \(\Ps(n, \beta) \ge 1 - \varepsilon\) and
  \(\Pt(2n, \beta) \ge 1 - \varepsilon\), then
  \(\Ps(2n, 2\beta) \ge 1 - 5\varepsilon^2\) and \(\Pt(4n, 2\beta) \ge 1 - 5\varepsilon^2\).
\end{lemma}
\begin{proof}
  \leanok
  This is Lemma~\ref{lem:doubling-coefficient} with \(C = 5\), since \(G\) is increasing and
  \(G(1/100) \le 5\). Directly: combining the first coupling identity with \(\Pt(n, \beta) \le 1\) gives
  \(\Ps(n, 2\beta) \ge (1-\varepsilon)^2 (1-\varepsilon)\), and spatial squaring
  (Lemma~\ref{lem:purity-squaring}) yields \(\Ps(2n, 2\beta) \ge 1 - 5\varepsilon^2\). Combining
  the second coupling identity with \(\Ps(2n, \beta) \le 1\) gives
  \(\Pt(4n, \beta) \ge (1-\varepsilon)^2 \Ps(2n, 2\beta) \ge (1-\varepsilon)^3\) since
  \(5\varepsilon^2 \le \varepsilon\), and thermal squaring yields
  \(\Pt(4n, 2\beta) \ge 1 - 5\varepsilon^2\).
\end{proof}

\begin{definition}[Defect sequence]
  \label{def:defect}
  \lean{CertifiedGaps.defect, CertifiedGaps.defect_succ, CertifiedGaps.defect_le}
  \leanok
  For \(\varepsilon_0 \in \R\), define \(\varepsilon_j = \tfrac15 (5\varepsilon_0)^{2^j}\). It
  satisfies \(\varepsilon_{j+1} = 5\varepsilon_j^2\), and \(\varepsilon_j \le 1/100\) whenever
  \(0 \le \varepsilon_0 \le 1/100\).
\end{definition}

\begin{lemma}[Induction on scale]
  \label{lem:induction-on-scale}
  \lean{CertifiedGaps.purities_ge_defect}
  \leanok
  \uses{lem:doubling, def:defect}
  Assume the spatial moment inequalities above an even \(n_0\), \(\beta_0 > 0\),
  \(0 < \varepsilon_0 \le 1/100\), and seeds \(\Ps(n_0, \beta_0) \ge 1 - \varepsilon_0\),
  \(\Pt(2n_0, \beta_0) \ge 1 - \varepsilon_0\). Then for every \(j\), with \(n_j = n_0 2^j\) and
  \(\beta_j = \beta_0 2^j\),
  \(\Ps(n_j, \beta_j) \ge 1 - \varepsilon_j\) and \(\Pt(2n_j, \beta_j) \ge 1 - \varepsilon_j\).
\end{lemma}
\begin{proof}
  \leanok
  Induct on \(j\), applying Lemma~\ref{lem:doubling} at \((n_j, \beta_j, \varepsilon_j)\).
\end{proof}

\begin{lemma}[Dyadic windows]
  \label{lem:dyadic-window}
  \lean{CertifiedGaps.exists_dyadic_window}
  \leanok
  If \(0 < n_0 \le L\), then \(n_0 2^j \le L \le 2 n_0 2^j\) for some \(j \in \N\).
\end{lemma}
\begin{proof}
  \leanok
  Take \(j = \lfloor \log_2 (L / n_0) \rfloor\) using integer division.
\end{proof}

\begin{theorem}[Two-purity criterion with a coefficient]
  \label{thm:gap-of-coefficient}
  \lean{CertifiedGaps.gapGT_of_coefficient, CertifiedGaps.coefficientDefect,
    CertifiedGaps.purities_ge_coefficientDefect}
  \leanok
  \uses{lem:doubling-coefficient, lem:dyadic-window, lem:purity-to-gap, def:moment-bounds,
    def:gapgt}
  This is Proposition~4.3 of~\cite{OpenAIHaldane}. Assume the spatial moment inequalities above an
  even \(n_0 > 0\), and let \(\beta_0 > 0\), \(0 < u_0 < 1/6\), \(G(u_0) \le C\),
  \(r = Cu_0 < 1\) and \(C(1 - 3u_0) > 3\). If \(\Ps(n_0, \beta_0) \ge 1 - u_0\) and
  \(\Pt(2n_0, \beta_0) \ge 1 - u_0\), then \(\GapGT\bigl(E_L, \log(1/r)/\beta_0\bigr)\) for every
  even \(L \ge n_0\).
\end{theorem}
\begin{proof}
  \leanok
  With \(u_j = r^{2^j}/C\), so that \(u_{j+1} = Cu_j^2 \le u_j \le u_0\), induction on \(j\) by
  Lemma~\ref{lem:doubling-coefficient} gives \(\Ps(n_j, \beta_j), \Pt(2n_j, \beta_j) \ge 1 - u_j\).
  Interpolation across the dyadic window of \(L\) gives \(\Pt(L, \beta_j) \ge (1-u_j)^3\), and
  Lemma~\ref{lem:purity-to-gap} bounds \(e^{-\beta_j \Delta_L}\) by
  \(3u_j / (1 - 3u_j) \le (3/C(1-3u_0))\, r^{2^j} < r^{2^j}\).
\end{proof}

\begin{theorem}[Two-purity criterion]
  \label{thm:gap-of-seeds}
  \lean{CertifiedGaps.gapGT_of_seeds}
  \leanok
  \uses{thm:gap-of-coefficient, lem:induction-on-scale}
  Assume the spatial moment inequalities above an even \(n_0 > 0\). Let \(\beta_0 > 0\) and
  \(0 < \varepsilon_0 \le 1/100\), with \(\Ps(n_0, \beta_0) \ge 1 - \varepsilon_0\) and
  \(\Pt(2n_0, \beta_0) \ge 1 - \varepsilon_0\). Then for every even \(L \ge n_0\),
  \[
    \GapGT\Bigl(E_L,\ \frac{\log\bigl(1/(5\varepsilon_0)\bigr)}{\beta_0}\Bigr).
  \]
\end{theorem}
\begin{proof}
  \leanok
  This is Theorem~\ref{thm:gap-of-coefficient} with \(u_0 = \varepsilon_0\) and \(C = 5\), so that
  \(r = 5\varepsilon_0\). In detail: choose the dyadic window \(n_j \le L \le 2n_j\) from Lemma~\ref{lem:dyadic-window} and set
  \(\beta = \beta_j\) and \(\varepsilon = \varepsilon_j\). To interpolate across the window, let
  \(r = L/(2n_j) \in (0, 1]\); inequality (3) of Definition~\ref{def:moment-bounds} at \(\beta\)
  and at \(2\beta\) gives \(Z_L(\beta)^2 \le Z_{n_j}(\beta)^{4r}\) and
  \(Z_{2n_j}(2\beta)^r \le Z_L(2\beta)\), so
  \[
    \Pt(L, \beta) \ge x^r \ge x, \qquad
    x = \frac{Z_{2n_j}(2\beta)}{Z_{n_j}(\beta)^4} = \Pt(2n_j, \beta)\, \Ps(n_j, \beta)^2
      \in \bigl[(1-\varepsilon)^3, 1\bigr],
  \]
  and hence \(\Pt(L, \beta) \ge 1 - 3\varepsilon\) by Lemma~\ref{lem:induction-on-scale}.
  Applying Lemma~\ref{lem:purity-to-gap} at a ground state \(m_0\) then bounds
  \(e^{-\beta (E_L(m) - E_L(m_0))}\) by \((1 - \Pt)/\Pt < (5\varepsilon_0)^{2^j}\). Taking
  logarithms and dividing by \(\beta = 2^j \beta_0\) gives
  \(E_L(m) - E_L(m_0) > \log(1/(5\varepsilon_0)) / \beta_0\).
\end{proof}

\begin{corollary}[The criterion with purities \(\ge 0.99\)]
  \label{cor:gap-of-purities}
  \lean{CertifiedGaps.gapGT_of_purities}
  \leanok
  \uses{thm:gap-of-seeds, prop:moment-form-bounds}
  Assume the spatial moment form above an even \(n_0 > 0\) and \(\Ps(n_0, \beta_0),
  \Pt(2n_0, \beta_0) \ge 0.99\). Then \(\GapGT(E_L, \log 20 / \beta_0)\) holds for every even
  \(L \ge n_0\).
\end{corollary}
\begin{proof}
  \leanok
  Apply Theorem~\ref{thm:gap-of-seeds} with \(\varepsilon_0 = 1/100\), using
  Proposition~\ref{prop:moment-form-bounds}.
\end{proof}
